% !TeX program = XeLaTeX
% This command controls the TeX engine you use: 
%				pdfLaTeX		-	Standard and fastest engine, but limited support for modern and unicode fonts.
%				XeLaTeX			-	Supports modern and unicode font.
%				LuaLaTeX		-	Most modern engine, slowest but supports all fonts and languages, additionally allows Lua code
% It is reccomended to stick to pdfLaTeX unless you really need to switch.
%%%%%%%%%%%%%%%%%%%%%%%%%%%%%%%%%%%%%%%%%%%%%%%%%%%%%%%%%%%%%%%%%%%%%%%%%%%%%%%%%%%%%%%%%%%%%%%%%%
% Thesis template
%%%%%%%%%%%%%%%%%%%%%%%%%%%%%%%%%%%%%%%%%%%%%%%%%%%%%%%%%%%%%%%%%%%%%%%%%%%%%%%%%%%%%%%%%%%%%%%%%%
% Version 1.1.2 by G. H. Allison

% This is an accessibility command that allow the document to auto tag.
% Delete this if it is causing problems
% If you are using Overleaf this will cause problems
\DocumentMetadata{
	lang        = en,
	pdfversion  = 2.0,
	pdfstandard = ua-2,
	pdfstandard = a-4f, %or a-4
	testphase   = 
	{phase-III,
		title,
		table,
		math,
		firstaid}  
}

% How to use this template:
% Everywhere you will need to edit is marked with XXXXXX in a comment, these comments also tell you 
% what needs to go where. Do not touch anything that is not commented if you do not understand it

% A note on fonts:
% When you install latex it will default to using pdflatex as the compiler. This dosn't support custom fonts. If you wish to use them, you must set the compiler to either XeLaTeX or LuaLaTeX.

% \RequirePackage{fix-cm}		% This provides font scaling for computer modern. The default LaTeX font. Comment this out if you are using another font.

\documentclass[12pt]{thesis} % tell the compiler to use the template

% XXXXXX load any additional packages, don't worry too much about this now, as you go along you will likely need to use some special packages for particular problems like full page tables etc.. This is where you tell the compiler what packages it should load before it tried to build your document.

\input{preamble.tex} % provides preamble

\title{Remember, Record, Resist:}     %note \\[1ex] is a line break in the title for nicer spacing if you wish to force it.

% Otherwise the title should wrap and adjust size automatically.
% If your title still overruns you should probably consider a shorter title

\subtitle{A Self-Contained QR-Code Archival System for Long-Term Information Preservation}

%%%%%%%%%%%%%%%%%%%%%%%%%%%%%%%%%%%%%%%%%%%%%%%%%%%%%%%%%%%%%%%%%%%%%%%%%%%%%%%%%%%%%%%%%%%%
%note if you want to apply a custom logo it can be written below here
% The default logo is defined as logo.pdf
% You can change the logo by placing a new logo.pdf file or using the customlogo commands below.
% the command takes multiple inputs seperated by a `,` and should size automatically.
% Multiple logos are defined by height rather than width.
%%%%%%%%%%%%%%%%%%%%%%%%%%%%%%%%%%%%%%%%%%%%%%%%%%%%%%%%%%%%%%%%%%%%%%%%%%%%%%%%%%%%%%%%%%%%

%\setcustomlogo{logo2.pdf, logo.pdf}	%XXXXXXX custom logos
%\setcustomlogo{logo2.pdf}

\uniname{PLACEHOLDER: University}  %XXXXXXX your university name

\author{A.K.A LYRA PHASMA}                           %XXXXXXX your name

\super{PLACEHOLDER: Supervisor} %note this command changes the supervisor listing based on number included

%\renewcommand{\submittedtext}{change the default text here if needed} 
\award{Bachelor of Science in Computer Science}                     %XXXXXXX the degree name

\faculty{College of Engineering and Information Technology}                  %XXXXXXX The name of the faculty for your school
\school{Department of Information Technology} %XXXXXXX the full name of the school

\degreedate{2027}                          %XXXXXXX the degree date
%end the preamble and start the document

\begin{document}
% \DeclareDatamodelFields[type=field,datatype=date]{eventdate}
% \DeclareDatamodelFields[type=field,datatype=date]{urldate}

%this baselineskip gives sufficient line spacing for an examiner to easily
%markup the thesis with comments

%XXXXXX when the compiler hits an \include{file_name} or \subfile{file_name} command it switches to that file goes all the way through it and comes back to this one. When you are working on single chapters you can comment out the other chapters so the document builds faster

\frontmatter

\begin{romanpages}          % start roman page numbering
\maketitle                  % create a title page from the preamble info
\subfile{1_thesis/statement_of_originality}        % include the statement_of_originality.tex file
\subfile{1_thesis/dedication}        % include the dedication.tex file
\subfile{1_thesis/acknowledgements}   % include the acknowledgements.tex file
\subfile{1_thesis/abstract}          % include the abstract.tex
\tableofcontents            % generate and include a table of contents
\listoffigures              % generate and include a list of figures
\listoftables				% generate and include a list of tables
\subfile{1_thesis/preface}        % include the preface.tex file
\end{romanpages}            % end roman page numbering

%XXXXXXXnow include the files of latex for each of the etc
% \subfile{chp1/doc1}
% \subfile{chp2/doc2}

\mainmatter

\subfile{1_thesis/chapters/1_introduction}

\subfile{1_thesis/chapters/2_rrl}

\subfile{1_thesis/chapters/3_methodology}

\appendix
% \part*{Appendix: Corpus of Arcane Texts}
% \addcontentsline{toc}{part}{Appendix: Corpus of Arcane Texts}
% \setcounter{chapter}{0}
% \renewcommand{\thechapter}{\arabic{chapter}}

% \subfile{math/main}

% \subfile{qr/main}

% \subfile{conclusions}

\printbibliography

%now enable appendix numbering format and include any appendices
\appendix

\backmatter

\printglossary

\end{document}